\documentclass[11pt]{article}
\newcommand{\doctitle}{Biomarker--Treatment Correlation During Carryover:
  A Modeling Artifact in the pmsimstats
  Data-Generating Process}
\newcommand{\shorttitle}{Carryover Correlation Artifact}
\newcommand{\docdate}{2026-03-19}
% pmsimstats-ng standard document preamble
% Usage: % pmsimstats-ng standard document preamble
% Usage: % pmsimstats-ng standard document preamble
% Usage: \input{pmsimstats-preamble} after \documentclass[11pt]{article}
%
% Documents must define before \input:
%   \newcommand{\doctitle}{...}       Full document title
%   \newcommand{\shorttitle}{...}     Short title for header
%   \newcommand{\docdate}{YYYY-MM-DD} Document date

\usepackage[margin=1in]{geometry}
\usepackage{amsmath,amssymb}
\usepackage[T1]{fontenc}
\usepackage{lmodern}
\usepackage{microtype}
\usepackage{graphicx}
\graphicspath{{figures/}}
\usepackage{booktabs}
\usepackage{longtable}
\usepackage{array}
\usepackage{hyperref}
\usepackage{xcolor}
\usepackage{fancyhdr}
\usepackage{parskip}
\usepackage{caption}
\usepackage{enumitem}
\usepackage{verbatim}
\usepackage{listings}

\lstset{
  language=R,
  basicstyle=\small\ttfamily,
  keywordstyle=\color{blue!70!black},
  commentstyle=\color{green!50!black},
  stringstyle=\color{red!60!black},
  breaklines=true,
  frame=single,
  framerule=0.5pt,
  rulecolor=\color{gray!40},
  backgroundcolor=\color{gray!5},
  xleftmargin=1em,
  framexleftmargin=1em,
  columns=flexible
}

\hypersetup{
  colorlinks=true,
  linkcolor=blue!60!black,
  citecolor=green!50!black,
  urlcolor=blue!60!black
}

\pagestyle{fancy}
\fancyhf{}
\lhead{\footnotesize pmsimstats-ng Technical Report}
\rhead{\footnotesize \shorttitle}
\rfoot{\footnotesize \thepage}
\renewcommand{\headrulewidth}{0.4pt}

\title{\doctitle}
\author{pmsimstats team}
\date{\docdate}
 after \documentclass[11pt]{article}
%
% Documents must define before \input:
%   \newcommand{\doctitle}{...}       Full document title
%   \newcommand{\shorttitle}{...}     Short title for header
%   \newcommand{\docdate}{YYYY-MM-DD} Document date

\usepackage[margin=1in]{geometry}
\usepackage{amsmath,amssymb}
\usepackage[T1]{fontenc}
\usepackage{lmodern}
\usepackage{microtype}
\usepackage{graphicx}
\graphicspath{{figures/}}
\usepackage{booktabs}
\usepackage{longtable}
\usepackage{array}
\usepackage{hyperref}
\usepackage{xcolor}
\usepackage{fancyhdr}
\usepackage{parskip}
\usepackage{caption}
\usepackage{enumitem}
\usepackage{verbatim}
\usepackage{listings}

\lstset{
  language=R,
  basicstyle=\small\ttfamily,
  keywordstyle=\color{blue!70!black},
  commentstyle=\color{green!50!black},
  stringstyle=\color{red!60!black},
  breaklines=true,
  frame=single,
  framerule=0.5pt,
  rulecolor=\color{gray!40},
  backgroundcolor=\color{gray!5},
  xleftmargin=1em,
  framexleftmargin=1em,
  columns=flexible
}

\hypersetup{
  colorlinks=true,
  linkcolor=blue!60!black,
  citecolor=green!50!black,
  urlcolor=blue!60!black
}

\pagestyle{fancy}
\fancyhf{}
\lhead{\footnotesize pmsimstats-ng Technical Report}
\rhead{\footnotesize \shorttitle}
\rfoot{\footnotesize \thepage}
\renewcommand{\headrulewidth}{0.4pt}

\title{\doctitle}
\author{pmsimstats team}
\date{\docdate}
 after \documentclass[11pt]{article}
%
% Documents must define before \input:
%   \newcommand{\doctitle}{...}       Full document title
%   \newcommand{\shorttitle}{...}     Short title for header
%   \newcommand{\docdate}{YYYY-MM-DD} Document date

\usepackage[margin=1in]{geometry}
\usepackage{amsmath,amssymb}
\usepackage[T1]{fontenc}
\usepackage{lmodern}
\usepackage{microtype}
\usepackage{graphicx}
\graphicspath{{figures/}}
\usepackage{booktabs}
\usepackage{longtable}
\usepackage{array}
\usepackage{hyperref}
\usepackage{xcolor}
\usepackage{fancyhdr}
\usepackage{parskip}
\usepackage{caption}
\usepackage{enumitem}
\usepackage{verbatim}
\usepackage{listings}

\lstset{
  language=R,
  basicstyle=\small\ttfamily,
  keywordstyle=\color{blue!70!black},
  commentstyle=\color{green!50!black},
  stringstyle=\color{red!60!black},
  breaklines=true,
  frame=single,
  framerule=0.5pt,
  rulecolor=\color{gray!40},
  backgroundcolor=\color{gray!5},
  xleftmargin=1em,
  framexleftmargin=1em,
  columns=flexible
}

\hypersetup{
  colorlinks=true,
  linkcolor=blue!60!black,
  citecolor=green!50!black,
  urlcolor=blue!60!black
}

\pagestyle{fancy}
\fancyhf{}
\lhead{\footnotesize pmsimstats-ng Technical Report}
\rhead{\footnotesize \shorttitle}
\rfoot{\footnotesize \thepage}
\renewcommand{\headrulewidth}{0.4pt}

\title{\doctitle}
\author{pmsimstats team}
\date{\docdate}


\begin{document}
\maketitle
\thispagestyle{fancy}

\begin{abstract}
The \texttt{pmsimstats} package (Hendrickson et al., 2020)
simulates clinical trial data by drawing participant outcomes
from a multivariate normal distribution whose correlation
structure encodes a biomarker--treatment interaction. A
post-publication commit introduced a change to the rule
governing how the biomarker--biologic response (BM--BR)
correlation is assigned during off-drug periods with carryover
residual. This white paper demonstrates that the original
correlation rule creates a modeling artifact: it applies full
biomarker predictive power to timepoints where the drug effect
has decayed to near zero, injecting spurious biomarker-specific
variance into off-drug periods. This artifact, not the carryover
residual itself, is the primary driver of the power loss reported
in the publication. This document traces the mathematics of the artifact,
evaluate both the original and revised correlation rules against
clinical and statistical criteria, and propose a principled
resolution.
\end{abstract}

\bigskip
\noindent\fbox{\parbox{\dimexpr\textwidth-2\fboxsep-2\fboxrule}{%
\textbf{Architecture scope.} The correlation artifact documented
here is specific to Architecture~B (MVN differential correlation),
in which the biomarker-BR correlation is assigned as a step
function of treatment status. Under Architecture~A (direct mean
moderation), the biomarker-treatment interaction is additive in
the mean structure and does not depend on the correlation
assignment logic. For a comparison of both DGP architectures, see
\texttt{02-dgp-mean-moderation-vs-mvn.tex}. For implementation
details of Architecture~A, see
\texttt{19-mean-moderation-implementation-notes.tex}.}}
\bigskip

\tableofcontents
\newpage

% =================================================================
\section{Introduction}
% =================================================================

The central finding of Hendrickson et al.\ (2020) is that the
N-of-1 trial design, while offering superior power under ideal
conditions, is acutely sensitive to carryover effects: even a
half-life of 0.1 weeks causes power to drop from 0.95 to 0.18
at $N = 35$. This finding has direct implications for trial
design decisions in precision medicine.

A post-publication commit (\texttt{8609f12}) by co-author Ron
Thomas modified the rule governing how the biomarker--biologic
response (BM--BR) correlation is assigned during off-drug
periods. Under the revised rule, the power loss from carryover
is eliminated entirely. The question of which rule is correct
is not a coding error to be fixed but a substantive modeling
choice with different implications for the data-generating
process (DGP).

This white paper provides a mathematical analysis of why the
original rule produces the power drop, whether the drop
reflects a genuine statistical phenomenon or a modeling
artifact, and what the implications are for the published
results.

% =================================================================
\section{The Data-Generating Process}
% =================================================================

\subsection{Multivariate Normal Framework}

Each simulated trial draws $N$ participants from a multivariate
normal distribution:
\[
\mathbf{X}_i \sim \mathcal{N}(\boldsymbol{\mu},\, \Sigma)
\]
where $\mathbf{X}_i$ is a vector of length $2 + 3K$ ($K$ =
number of timepoints), containing the biomarker (BM), baseline
symptoms (BL), and three response factors (TV, ER, BR) at each
timepoint. The observed outcome is:
\[
S_{i,t} = \text{BL}_i -
  (\text{TV}_{i,t} + \text{ER}_{i,t} + \text{BR}_{i,t})
\]

The biomarker--treatment interaction is encoded not in the
mean vector $\boldsymbol{\mu}$ but in the covariance matrix
$\Sigma$. Specifically, the correlation between the biomarker
and the BR factor at timepoint $t$ determines the extent to
which participants with different biomarker values have
different drug responses.

\subsection{Mean Vector Construction}

The BR mean at timepoint $t$ is computed from the modified
Gompertz function evaluated at the cumulative time on drug
(tod):
\[
\mu_{\text{BR},t}^{\text{raw}} =
  f_G(\text{tod}_t;\, \text{max}, \text{disp}, \text{rate})
\]
where $f_G(0) = 0$ by construction. When a participant is off
drug, $\text{tod}_t = 0$, so
$\mu_{\text{BR},t}^{\text{raw}} = 0$.

Carryover is applied sequentially:
\[
\mu_{\text{BR},t}^{\text{adj}} =
  \mu_{\text{BR},t}^{\text{raw}} +
  \mu_{\text{BR},t-1}^{\text{adj}} \cdot
  \left(\tfrac{1}{2}\right)^{s \cdot t_{sd} / t_{1/2}}
\]
where $t_{sd}$ is time since discontinuation, $t_{1/2}$ is
the carryover half-life, and $s$ is a scale factor (1 in the
original code, 2 after the Ron Thomas edit).

\subsection{Covariance Matrix Construction}

The covariance matrix is built from a correlation matrix
$R$ and a vector of standard deviations $\boldsymbol{\sigma}$:
\[
\Sigma = \text{diag}(\boldsymbol{\sigma}) \cdot R \cdot
  \text{diag}(\boldsymbol{\sigma})
\]

The correlation between BM and $\text{BR}_t$ is the critical
element. The two code versions assign this correlation
differently, as described in the following sections.

% =================================================================
\section{The Two Correlation Rules}
% =================================================================

\subsection{Original Rule (Commit \texttt{42ac030})}

The original code assigns the BM--BR correlation using a
simple threshold test on the adjusted BR mean:
\[
\rho(\text{BM}, \text{BR}_t) =
\begin{cases}
c_{bm} & \text{if } \mu_{\text{BR},t}^{\text{adj}} \neq 0 \\
0 & \text{if } \mu_{\text{BR},t}^{\text{adj}} = 0
\end{cases}
\]

This is a step function: any nonzero adjusted BR mean, no
matter how small, triggers the full correlation $c_{bm}$.

In the R code:
\begin{verbatim}
if (means[which(n1 == labels)] != 0) {
  correlations[n1, 'bm'] <- modelparam$c.bm
  correlations['bm', n1] <- modelparam$c.bm
}
\end{verbatim}

\subsection{Revised Rule (Commit \texttt{8609f12})}

The Ron Thomas revision introduces a three-way assignment with
proportional scaling during carryover periods:
\[
\rho(\text{BM}, \text{BR}_t) =
\begin{cases}
0 & \text{if } t = 1 \text{ (first timepoint)} \\[4pt]
c_{bm} & \text{if on drug (raw BR} > 0\text{)} \\[4pt]
\displaystyle\frac{\mu_{\text{BR},t}^{\text{adj}}}
  {\mu_{\text{BR},t-1}^{\text{adj}}} \cdot c_{bm}
  & \text{if off drug, raw BR} = 0,\,
    \mu_{\text{BR},t}^{\text{adj}} > 0 \\[4pt]
0 & \text{if off drug, }
  \mu_{\text{BR},t}^{\text{adj}} = 0
\end{cases}
\]

In the R code:
\begin{verbatim}
if (p > 1) {
  correlations["bm", n1] <- correlations[n1, "bm"] <-
    ifelse(brtest[p],
           ifelse(brmeans[p] == 0, 0,
                  (mm1 / mm0) * modelparam$c.bm),
           modelparam$c.bm)
}
\end{verbatim}

where \texttt{brtest[p]} is \texttt{TRUE} when the raw
(pre-carryover) BR mean is zero at timepoint $p$.

% =================================================================
\section{The Artifact: Conditional Expectations Under the
         Original Rule}
% =================================================================

\subsection{Conditional Mean of BR Given Biomarker}

In a multivariate normal distribution, the conditional
expectation of $\text{BR}_t$ given a specific biomarker
value $bm$ is:
\begin{equation}
E[\text{BR}_t \mid \text{BM} = bm] =
  \mu_{\text{BR},t} +
  \frac{\rho \cdot \sigma_{\text{BR}}}{\sigma_{\text{BM}}}
  \cdot (bm - \mu_{\text{BM}})
\label{eq:conditional}
\end{equation}

The key observation is that this conditional expectation
depends on $\rho$ and $\sigma_{\text{BR}}$ but \emph{not}
on $\mu_{\text{BR},t}$ (which only shifts the intercept).
When $\rho = c_{bm} = 0.6$ and $\sigma_{\text{BR}} = 8$,
$\sigma_{\text{BM}} = 15.36$:
\[
E[\text{BR}_t \mid \text{BM} = bm] =
  \mu_{\text{BR},t} + 0.3125 \cdot (bm - 124.3)
\]

The slope $0.3125$ is independent of $\mu_{\text{BR},t}$.

\subsection{Numerical Example: N-of-1 Path A at BD3}

Consider timepoint BD3 in the N-of-1 design (Path A), where the
participant has been off drug for 1 week with carryover
half-life $t_{1/2} = 0.1$ weeks:

\begin{itemize}[nosep]
\item Raw BR mean: $\mu_{\text{BR},5}^{\text{raw}} = 0$
  (off drug, tod $= 0$)
\item Adjusted BR mean: $\mu_{\text{BR},5}^{\text{adj}} = 0.01$
  (negligible carryover residual: $10.19 \times 2^{-10}$)
\item BR standard deviation: $\sigma_{\text{BR}} = 8$
  (unchanged regardless of mean)
\end{itemize}

Under the original rule, $\rho = 0.6$ because
$\mu_{\text{BR},5}^{\text{adj}} \neq 0$. Applying
Equation~\ref{eq:conditional}:

\begin{table}[h]
\centering
\caption{Conditional $E[\text{BR}_5 \mid \text{BM}]$ at
  BD3 under the original correlation rule}
\begin{tabular}{lrrr}
\toprule
Biomarker value & BM deviation & $E[\text{BR}]$
  (original) & $E[\text{BR}]$ (revised) \\
\midrule
109 ($-1$ SD) & $-15.4$ & $-4.8$ & $+0.01$ \\
124 (mean) & $0$ & $+0.01$ & $+0.01$ \\
140 ($+1$ SD) & $+15.7$ & $+4.9$ & $+0.01$ \\
\bottomrule
\end{tabular}
\label{tab:conditional}
\end{table}

Under the original rule, the conditional BR spans a range of
nearly 10 points ($-4.8$ to $+4.9$) centered on a mean of
0.01. The biomarker predicts a large, clinically meaningful
difference in drug effect at a timepoint where the drug effect
has decayed to 0.01\% of its peak.

Under the revised rule, $\rho \approx 10^{-6}$, and the
conditional BR is 0.01 for all biomarker values.

\subsection{The Source of the Artifact}

The artifact arises from a disconnect between two elements of
the multivariate normal specification:

\begin{enumerate}
\item \textbf{The mean vector} reflects the drug effect
  decaying to near zero via carryover.
\item \textbf{The correlation matrix} maintains full
  biomarker--BR correlation ($\rho = 0.6$) as long as the
  adjusted mean is nonzero.
\item \textbf{The standard deviation vector} is fixed at
  $\sigma_{\text{BR}} = 8$ regardless of the mean.
\end{enumerate}

The combination of full correlation and fixed variance around
a near-zero mean produces draws where some participants (high
biomarker) show substantial positive BR and others (low
biomarker) show substantial negative BR, despite the mean
drug effect being negligible. The negative BR values are
particularly problematic: they imply that the drug is
\emph{harming} low-biomarker participants during the off-drug
period, which has no clinical basis.

\subsection{Why This Destroys Power}

The LME analysis model for crossover designs is:
\[
S_{i,t} = \beta_0 + \beta_1 \cdot bm_i +
  \beta_2 \cdot D_{bc,t} + \beta_3 \cdot t +
  \beta_4 \cdot bm_i \cdot D_{bc,t} +
  u_i + \varepsilon_{i,t}
\]
where $D_{bc}$ is the drug indicator and $\beta_4$ is the
biomarker--treatment interaction of interest.

The spurious BM--BR correlation during off-drug periods
introduces structured noise into the $bm \cdot D_{bc}$
term. When the participant is off drug ($D_{bc} = 0$ or
near zero), the model expects no biomarker--drug interaction.
But the correlated BR draws create biomarker-specific variation
that inflates the residual variance of the interaction
estimate, increasing the standard error of $\hat{\beta}_4$
and reducing the probability of rejecting $H_0: \beta_4 = 0$.

The effect is largest for designs with many off-drug timepoints
(N-of-1 has 4 out of 8) and smallest for designs with no
off-drug periods (OL is unaffected because all timepoints are
on drug).

% =================================================================
\section{Clinical Evaluation}
% =================================================================

\subsection{Is the Original Rule Clinically Reasonable?}

The original rule implies the following clinical scenario:

\begin{quote}
One week after discontinuing a drug whose pharmacologic
effect has decayed to 0.01\% of its peak value, a patient's
baseline blood pressure still perfectly predicts whether the
patient is ``responding'' (positive BR) or ``not responding''
(negative BR) to the drug.
\end{quote}

This is not clinically plausible. The biomarker--treatment
interaction is mediated through the drug's mechanism of action.
When the drug effect has decayed to negligible levels, the
biomarker's predictive value for the drug effect should also
decay. The step-function behavior of the original rule (full
predictive power until the last molecule clears, then
instantly zero) does not correspond to any known pharmacologic
mechanism.

Furthermore, the negative conditional BR values for
low-biomarker participants are clinically incoherent. They
imply that discontinuing the drug \emph{worsens} symptoms
beyond the natural trajectory, and that this worsening is
specific to low-biomarker individuals. No such rebound effect
was hypothesized in the original publication.

\subsection{Is the Revised Rule Clinically Reasonable?}

The revised rule implies:

\begin{quote}
The biomarker's predictive power for the drug effect decays
in exact proportion to the drug effect itself. When the drug
effect is at 50\% of peak, the biomarker--BR correlation is
$0.5 \cdot c_{bm}$. When the drug effect is at 0.01\%,
the correlation is $0.0001 \cdot c_{bm}$.
\end{quote}

This is more clinically reasonable but may be overly
optimistic. It assumes perfect proportional tracking between
the drug effect magnitude and the biomarker's predictive
value. In practice, the relationship might be non-linear:
the biomarker might retain some predictive value even as
the drug effect wanes (e.g., if the biomarker indexes a
slow-to-reverse physiological state that mediates the drug
response).

However, the practical difference between the revised rule
and any plausible intermediate model is negligible at the
carryover half-lives tested ($t_{1/2} = 0.1$ and $0.2$
weeks). At these values, the carryover residual at the next
assessment timepoint (1--4 weeks later) is less than 0.1\%
of the peak BR, and any correlation proportional to the
residual would be effectively zero.

\subsection{Summary of Clinical Assessment}

\begin{table}[h]
\centering
\caption{Clinical plausibility of the two correlation rules}
\begin{tabular}{lcc}
\toprule
Criterion & Original & Revised \\
\midrule
Drug effect near zero $\Rightarrow$ BM predictive?
  & Yes (full) & No ($\approx 0$) \\
Negative BR during off-drug periods?
  & Yes (for low BM) & No \\
Proportional decay of predictive power?
  & No (step function) & Yes (exact) \\
Clinically plausible?
  & \textcolor{red}{No} & Approximately \\
\bottomrule
\end{tabular}
\end{table}

% =================================================================
\section{Statistical Evaluation}
% =================================================================

\subsection{Covariance Matrix Coherence}

In a well-specified multivariate normal model, the correlation
structure should be consistent with the mean structure. When
$\mu_{\text{BR},t} \approx 0$, applying a large correlation
with BM creates a conditional distribution where most of the
conditional variation is biomarker-driven rather than
drug-driven. The conditional variance of BR given BM is:
\[
\text{Var}(\text{BR}_t \mid \text{BM}) =
  \sigma_{\text{BR}}^2 (1 - \rho^2) =
  64 \cdot (1 - 0.36) = 40.96
\]

The conditional standard deviation is $\sqrt{40.96} = 6.4$,
which is enormous relative to the mean of 0.01. The
coefficient of variation is 640, indicating a distribution
centered at effectively zero but with substantial spread
driven by the biomarker conditioning. This is statistically
valid (the MVN can produce such draws) but represents a
pathological parameterization where the correlation structure
dominates the mean structure at off-drug timepoints.

\subsection{Impact on the Covariance Matrix Positive
            Definiteness}

The original rule can create situations where the correlation
between BM and BR changes abruptly between adjacent timepoints
(e.g., from 0.6 at BD2 to 0.6 at BD3 to 0.6 at BD4, despite
BD3--BD4 being off drug). This does not violate positive
definiteness per se, because the same correlation is applied
to all off-drug timepoints with any residual. However, the
step function creates discontinuities in the correlation
structure that increase the likelihood of requiring the
\texttt{make.positive.definite()} correction, which in turn
perturbs the intended correlation structure.

\subsection{Effect Size Bias}

The spurious BM--BR correlation during off-drug periods does
not bias the point estimate of $\beta_4$ (the interaction
coefficient) in expectation, because the noise is symmetric
around zero. However, it inflates the standard error of
$\hat{\beta}_4$ by introducing structured residual variance
that the random effects cannot absorb. The practical
consequence is not bias but loss of precision, hence loss
of power.

\subsection{Type I Error}

At $c_{bm} = 0$ (no true interaction), neither rule produces
spurious correlations, because both rules multiply $c_{bm}$
by some factor. When $c_{bm} = 0$, any multiplicative
scaling yields zero. Consequently, the Type I error rate is
approximately 0.05 under both rules, consistent with the
publication's findings and the present simulation results.

% =================================================================
\section{Implications for Published Results}
% =================================================================

\subsection{Which Results Are Affected?}

Only results involving nonzero carryover ($t_{1/2} > 0$) and
designs with off-drug timepoints are affected. Specifically:

\begin{itemize}[nosep]
\item \textbf{Figure 4A} (carryover $= 0$): \textbf{Not
  affected.} When carryover is zero, there are no carryover
  residuals, and both rules assign $\rho = 0$ to off-drug
  timepoints with $\mu_{\text{BR}} = 0$.
\item \textbf{Figure 4B} (carryover $> 0$): \textbf{Affected.}
  The power drop for N-of-1 and OL+BDC is driven by the
  correlation artifact, not by the carryover residual itself.
\item \textbf{OL design}: \textbf{Not affected.} All
  timepoints are on drug, so no off-drug correlation
  assignment occurs.
\end{itemize}

\subsection{Is the Paper's Central Conclusion Still Valid?}

The paper's central conclusion is that N-of-1 designs are
sensitive to carryover effects. Under the revised correlation
rule, this conclusion does not hold: N-of-1 power is robust
to carryover at the tested half-lives.

However, the paper's conclusion may still be qualitatively
valid for a different reason. The carryover \emph{mean}
residual at $t_{1/2} = 0.1$ weeks is negligible (0.01\%),
but at longer half-lives (e.g., $t_{1/2} = 2$ weeks), the
residual would be substantial:
\[
\mu_{\text{BR},t}^{\text{adj}} =
  10.19 \times (1/2)^{1/2} \approx 7.2
  \quad\text{at } t_{1/2} = 2, t_{sd} = 1
\]

At such half-lives, even the revised correlation rule would
assign a substantial BM--BR correlation ($\sim 0.7 \times
c_{bm}$), and the carryover would genuinely confound the
analysis by maintaining a biomarker-specific drug effect
during off-drug periods. The paper's finding would hold, but
at longer carryover half-lives than those tested.

\subsection{Recommendation}

The published results at carryover half-lives of 0.1 and 0.2
weeks should be interpreted with caution. The dramatic power
loss reported is primarily an artifact of the correlation
assignment rule, not a genuine statistical consequence of
carryover at those half-lives.

Future simulations should either:

\begin{enumerate}
\item Use the revised (proportional) correlation rule and
  test carryover at longer, clinically realistic half-lives
  (e.g., 1--4 weeks), or
\item Explicitly model the correlation decay function as a
  separate parameter, allowing sensitivity analysis over
  different assumptions about whether carryover effects
  retain biomarker specificity.
\end{enumerate}

% =================================================================
\section{A Principled Correlation Decay Model}
% =================================================================

Neither the step function (original) nor exact proportional
scaling (revised) is fully satisfactory as a general model.
A more principled approach would parameterize the correlation
decay explicitly:

\[
\rho(\text{BM}, \text{BR}_t) =
\begin{cases}
c_{bm} & \text{if on drug} \\[4pt]
c_{bm} \cdot g(t_{sd};\, \boldsymbol{\theta}) &
  \text{if off drug}
\end{cases}
\]

where $g(t_{sd};\, \boldsymbol{\theta})$ is a decay function
with parameters $\boldsymbol{\theta}$ controlling the rate
at which biomarker predictive power decays after drug
discontinuation. Natural choices include:

\begin{itemize}
\item \textbf{Step function} (original):
  $g(t_{sd}) = \mathbf{1}[\mu_{\text{BR},t}^{\text{adj}} > 0]$
\item \textbf{Proportional} (revised):
  $g(t_{sd}) = \mu_{\text{BR},t}^{\text{adj}} /
  \mu_{\text{BR},t-1}^{\text{adj}}$
\item \textbf{Exponential decay}:
  $g(t_{sd}) = \exp(-\lambda \cdot t_{sd})$ with separate
  half-life $\lambda$ for the correlation
\item \textbf{Power law}:
  $g(t_{sd}) = (1 + t_{sd})^{-\alpha}$
\end{itemize}

The exponential decay model is particularly attractive because
it allows the correlation half-life to differ from the mean
effect half-life, which is pharmacologically plausible: the
biomarker's predictive power might decay faster or slower
than the drug effect depending on the mechanism of action.

This parameterization would allow sensitivity analysis over
$\lambda$ and separate the question of ``how long does the
drug effect last?'' from ``how long does the biomarker remain
predictive?''

% =================================================================
\section{Conclusions}
% =================================================================

\begin{enumerate}
\item The power drop reported in Hendrickson et al.\ (2020)
  Figure~4B at carryover half-lives of 0.1 and 0.2 weeks is
  primarily driven by a modeling artifact in the correlation
  assignment rule, not by the carryover residual in the mean
  structure.

\item The artifact arises because the original code applies
  full biomarker--BR correlation ($\rho = c_{bm}$) to
  off-drug timepoints where the adjusted BR mean is
  negligible, creating clinically implausible conditional
  expectations (participants with low biomarkers show drug
  ``harm'' during off-drug periods).

\item The Ron Thomas proportional scaling (commit
  \texttt{8609f12}) corrects this artifact by scaling the
  correlation proportional to the carryover residual. This
  is clinically and statistically more defensible than the
  step-function rule.

\item The paper's qualitative conclusion (that N-of-1
  designs are sensitive to carryover) may still hold at
  longer, more clinically realistic carryover half-lives,
  but the quantitative results at $t_{1/2} = 0.1$ and
  $0.2$ weeks should be interpreted with caution.

\item Future implementations should parameterize the
  correlation decay function independently from the mean
  decay function, allowing explicit sensitivity analysis
  over assumptions about whether carryover drug effects
  retain biomarker specificity.
\end{enumerate}

% =================================================================
\section*{Appendix: Commit Chronology and Code Diffs}
% =================================================================

The following commits modified the core DGP files
(\texttt{generateData.R} and \texttt{lme\_analysis.R}):

\begin{description}
\item[\texttt{42ac030} (Initial commit)]
  Published code. BM--BR correlation assigned as step
  function on adjusted BR mean. No scale factor on
  carryover decay.

\item[\texttt{8609f12} (Ron Thomas edits)]
  Three changes: (1)~added scale factor $s = 2$ on
  carryover decay, (2)~replaced step-function correlation
  rule with proportional scaling, (3)~added timepoint-1
  zero-correlation convention. Also added
  \texttt{full\_model\_out} option to
  \texttt{lme\_analysis()}.

\item[\texttt{f6ee86d} (Autocorrelation typo fix)]
  Fixed a symmetry bug in the cross-factor correlation
  matrix (\texttt{correlations[n2,n1]} was incorrectly
  written as \texttt{correlations[n1,n2]}). No impact on
  Figure~4 because \texttt{make.positive.definite()} had
  been implicitly correcting the asymmetry.

\item[\texttt{f70f86d} (Carryover term added)]
  Added carryover functionality to other parts of the
  codebase. No changes to the core \texttt{generateData()}
  or \texttt{lme\_analysis()} functions that affect
  Figure~4.

\item[\texttt{325314f} (Modified Db to include carryover)]
  Introduced continuous drug indicator \texttt{Dbc} in the
  analysis model. Contains two bugs preventing execution
  (typo in \texttt{op*carryover\_scalefactor} and incorrect
  coefficient name \texttt{bm:DbcTRUE}).

\item[\texttt{HEAD} (Current code)]
  Fixed version of \texttt{325314f}. Uses continuous
  \texttt{Dbc} with exponential decay in the analysis
  model. Combined with the DGP changes from
  \texttt{8609f12}, produces the lowest correlation with
  the published Figure~4 ($r = 0.613$).
\end{description}

\vfill
\noindent\rule{\textwidth}{0.4pt}
{\footnotesize
Rendered on 2026-03-25 at 18:49 PDT.\\
Source: \textasciitilde/prj/alz/10-pmsimstats-ng/pmsimstats-ng/docs/09-carryover-correlation-artifact.tex}

\end{document}
